\uuid{gmi6}
\exo7id{2735}
\titre{exo7 2735}
\auteur{tumpach}
\organisation{exo7}
\datecreate{2009-10-25}
\isIndication{false}
\isCorrection{false}
\chapitre{Déterminant, système linéaire}
\sousChapitre{Système linéaire, rang}
\cpge{mpsi-mp2i.calcul-matriciel-systemes.c, pcsi.calcul-matriciel-systemes-lineaires.c, ptsi.calcul-matriciel-systemes-lineaires.c, mpsi-mp2i.matrices-applications-lineaires.c, pcsi.matrices-applications-lineaires.d, ptsi.matrices-applications-lineaires.d, bcpst1.systemes-lineaires.a, bcpst1.systemes-lineaires.c}
\module{Algèbre}
\niveau{L2}
\difficulte{}

\contenu{
\texte{
Soit $a$ un nombre r\'eel. On \'etudie le syst\`eme lin\'eaire suivant~:
$$
\mathcal{S}_{a}~: \left\{\begin{array}{ccccccc} 
x & - & 2y & + & 3z & = & 2\\
x & + & 3y & - & 2z & = & 5\\
2x & - & y & + & az & = & 1
\end{array}\right.
$$
}
\begin{enumerate}
    \item \question{En fonction des valeurs du param\`etre $a$, d\'eterminer si le syst\`eme $\mathcal{S}_{a}$ peut~:
}
    \item \question{[(i)] n'admettre aucune solution ;}
    \item \question{[(ii)] admettre exactement une solution ;}
    \item \question{[(iii)] admettre une infinit\'e de solutions.}
\end{enumerate}
}
